\Askhsh{34448}{4}
\begin{ylh}{1}
    \item 2.2 Παρουσίαση Στατιστικών Δεδομένων
    \item 2.3 Μέτρα Θέσης και Διασποράς
\end{ylh}
Μελετήσαμε ένα δείγμα Ι.Χ. αυτοκινήτων που κυκλοφορούν στο κέντρο της Αθήνας ως προς τον αριθμό των επιβατών συμπεριλαμβανομένου και του οδηγού. Μερικά από τα αποτελέσματα φαίνονται στον παρακάτω πίνακα.

\begin{center}
\begin{mytblr}{}
{Αριθμός\\επιβατών $x_i$} & {Αριθμός\\αυτοκινήτων $\nu_i$} & {$f_i$} & {$f_i\%$} & {$N_i$} & {$F_i$} & {$F_i\%$} \\

1   & 50    & 0,125 &       &       &       &       \\

2   & 110   & 0,275 &       &       &       &       \\

3   &       & 0,3   &       &       &       &       \\

4   & 30    & 0,075 &       &       &       &       \\

5   & 90    &       &       &       &       &       \\

\textbf{ΣΥΝΟΛΑ} & $\nu=400$ & & &  &  &  \\

\end{mytblr}
\end{center}

\begin{alist}
    \item Να μεταφέρετε τον παραπάνω πίνακα στο τετράδιό σας και να τον συμπληρώσετε. \monades{7}
    \item Να υπολογίσετε την μέση τιμή και τη διάμεσο του δείγματος. \monades{7}
    \item Ο Γιάννης ισχυρίζεται ότι το δείγμα είναι ομοιογενές. Συμφωνείτε με τον ισχυρισμό του Γιάννη; Δικαιολογήστε πλήρως την απάντησή σας. \monades{11}
\end{alist}
(Δίνεται ότι: $\sqrt{7} = 2,65$)

